%% ============================================================
%%  example.tex — Full demonstration of mermaid-overleaf.sty
%%
%%  HOW TO USE ON OVERLEAF:
%%  1. In your Overleaf project, click "New File" → "From External URL"
%%  2. Paste: https://raw.githubusercontent.com/YOUR-USERNAME/mermaid-overleaf/main/mermaid-overleaf.sty
%%  3. Save the file as mermaid-overleaf.sty
%%  4. \usepackage{mermaid-overleaf} in your document
%% ============================================================
\documentclass[12pt, a4paper]{article}

%% Load the package — default theme is "ocean"
\usepackage[theme=ocean]{mermaid-overleaf}

%% Standard document packages
\usepackage[T1]{fontenc}
\usepackage[utf8]{inputenc}
\usepackage{lmodern}
\usepackage{microtype}
\usepackage{hyperref}
\usepackage{geometry}
\geometry{margin=2.5cm}

\hypersetup{
  colorlinks=true,
  linkcolor=blue!60!black,
  urlcolor=blue!60!black,
}

\title{%
  \textbf{mermaid-overleaf}\\[0.4em]
  \large Mermaid Diagrams in Cloud Overleaf\\[0.2em]
  \normalsize Complete Usage Guide \& Examples
}
\author{mermaid-overleaf contributors}
\date{\today}

\begin{document}
\maketitle
\tableofcontents
\newpage

%% ============================================================
\section{Introduction}
%% ============================================================

The \textbf{mermaid-overleaf} package lets you write
\href{https://mermaid.js.org}{Mermaid} diagram source directly
inside your \LaTeX{} document and render it beautifully — even
on cloud Overleaf where shell-escape is not available.

\MermaidWorkflow

\bigskip

The package ships five built-in colour themes
(\texttt{ocean}, \texttt{forest}, \texttt{sunset},
\texttt{arctic}, \texttt{midnight}), auto-detects the diagram
type for its badge, and supports pre-rendered images seamlessly.

%% ============================================================
\section{Flowchart}
%% ============================================================

\subsection{Simple Decision Flow — \textit{ocean} theme}

\begin{mermaid}[type=flowchart, theme=ocean,
                caption={Order processing decision flow},
                label={fig:order-flow}]
flowchart TD
    A([Start]) --> B[Receive Order]
    B --> C{In Stock?}
    C -->|Yes| D[Process Payment]
    C -->|No|  E[Notify Customer]
    D --> F{Payment OK?}
    F -->|Yes| G[Ship Order]
    F -->|No|  H[Cancel Order]
    G --> I([End])
    H --> I
    E --> I
\end{mermaid}

\subsection{Left-to-Right Architecture — \textit{forest} theme}

\begin{mermaid}[type=flowchart, theme=forest,
                caption={Microservice architecture overview},
                label={fig:arch}]
flowchart LR
    subgraph Client
        A[Browser] --> B[Mobile App]
    end
    subgraph Gateway
        C[API Gateway]
        D[Rate Limiter]
    end
    subgraph Services
        E[Auth Service]
        F[Product Service]
        G[Order Service]
    end
    subgraph Data
        H[(Postgres)]
        I[(Redis Cache)]
        J[(S3 Storage)]
    end
    A & B --> C
    C --> D
    D --> E & F & G
    E --> H
    F --> H & I & J
    G --> H & I
\end{mermaid}

%% ============================================================
\section{Sequence Diagram}
%% ============================================================

\subsection{OAuth 2.0 Login Flow — \textit{midnight} theme}

\begin{mermaid}[type=sequence, theme=midnight,
                caption={OAuth 2.0 authorisation code flow},
                label={fig:oauth}]
sequenceDiagram
    actor U as User
    participant B as Browser
    participant A as App Server
    participant O as OAuth Provider
    participant R as Resource API

    U->>B: Click "Login with Google"
    B->>A: GET /auth/start
    A->>B: Redirect to OAuth URL
    B->>O: GET /authorize?client_id=...
    O->>U: Show consent screen
    U->>O: Approve
    O->>B: Redirect with code=XYZ
    B->>A: GET /callback?code=XYZ
    A->>O: POST /token (code + secret)
    O->>A: access_token + refresh_token
    A->>R: GET /api/user (Bearer token)
    R->>A: User profile JSON
    A->>B: Set session cookie
    B->>U: Dashboard loaded
\end{mermaid}

\subsection{Database Query — \textit{ocean} theme}

\begin{mermaid}[type=sequence, theme=ocean,
                caption={Cached database query sequence},
                label={fig:dbquery}]
sequenceDiagram
    participant App
    participant Cache as Redis Cache
    participant DB as PostgreSQL

    App->>Cache: GET user:1234
    alt Cache hit
        Cache-->>App: User JSON
    else Cache miss
        Cache-->>App: nil
        App->>DB: SELECT * FROM users WHERE id=1234
        DB-->>App: Row data
        App->>Cache: SET user:1234 (TTL 300s)
        Cache-->>App: OK
    end
    App->>App: Return user data
\end{mermaid}

%% ============================================================
\section{Gantt Chart}
%% ============================================================

\begin{mermaid}[type=gantt, theme=sunset,
                caption={Q1 2026 product roadmap},
                label={fig:gantt}]
gantt
    title Q1 2026 Product Roadmap
    dateFormat YYYY-MM-DD
    section Design
        UX Research           :done,    d1, 2026-01-01, 14d
        Wireframes            :done,    d2, after d1,   10d
        UI Design             :active,  d3, after d2,   14d
        Design Review         :         d4, after d3,   3d
    section Engineering
        Backend API           :done,    e1, 2026-01-05, 20d
        Frontend Components   :active,  e2, after d2,   21d
        Integration           :crit,    e3, after e1,   10d
        Performance Testing   :crit,    e4, after e3,   7d
    section Release
        Beta Launch           :milestone, 2026-02-28, 0d
        Bug Fixes             :          r1, 2026-03-01, 10d
        Public Release        :milestone, 2026-03-15, 0d
\end{mermaid}

%% ============================================================
\section{Class Diagram}
%% ============================================================

\begin{mermaid}[type=class, theme=arctic,
                caption={E-commerce domain class model},
                label={fig:classdiag}]
classDiagram
    class User {
        +int id
        +string email
        +string name
        +register()
        +login()
    }
    class Order {
        +int id
        +Date createdAt
        +string status
        +calculate_total() float
        +cancel()
    }
    class OrderItem {
        +int quantity
        +float unitPrice
        +subtotal() float
    }
    class Product {
        +int id
        +string name
        +float price
        +int stock
        +is_available() bool
    }
    class Payment {
        +string method
        +float amount
        +string status
        +process()
        +refund()
    }
    User "1" --> "0..*" Order : places
    Order "1" *-- "1..*" OrderItem : contains
    OrderItem "*" --> "1" Product : references
    Order "1" --> "1" Payment : paid via
\end{mermaid}

%% ============================================================
\section{State Diagram}
%% ============================================================

\begin{mermaid}[type=state, theme=forest,
                caption={Order lifecycle state machine},
                label={fig:state}]
stateDiagram-v2
    [*] --> Pending : Order placed

    Pending --> Confirmed : Payment OK
    Pending --> Cancelled : Payment failed

    Confirmed --> Preparing : Warehouse picks

    state Preparing {
        [*] --> Picking
        Picking --> Packing
        Packing --> ReadyToShip
        ReadyToShip --> [*]
    }

    Preparing --> Shipped : Courier collected
    Shipped --> Delivered : Customer received
    Shipped --> ReturnRequested : Customer disputes

    Delivered --> [*]
    ReturnRequested --> Refunded : Return approved
    ReturnRequested --> Delivered : Dispute rejected
    Refunded --> [*]
    Cancelled --> [*]
\end{mermaid}

%% ============================================================
\section{Entity-Relationship Diagram}
%% ============================================================

\begin{mermaid}[type=er, theme=midnight,
                caption={Blog platform ER diagram},
                label={fig:er}]
erDiagram
    USER {
        int     id     PK
        string  email
        string  name
        date    joined
    }
    POST {
        int     id     PK
        string  title
        text    body
        date    published
        int     author_id FK
    }
    TAG {
        int     id    PK
        string  name
    }
    COMMENT {
        int     id       PK
        text    content
        date    created
        int     post_id  FK
        int     user_id  FK
    }
    USER    ||--o{ POST    : writes
    POST    ||--o{ COMMENT : has
    USER    ||--o{ COMMENT : writes
    POST    }o--o{ TAG     : tagged-with
\end{mermaid}

%% ============================================================
\section{Timeline}
%% ============================================================

\begin{mermaid}[type=timeline, theme=ocean,
                caption={History of web development},
                label={fig:timeline}]
timeline
    title History of Web Development
    section 1990s
        1991 : World Wide Web born
             : First website by Tim Berners-Lee
        1994 : Netscape Navigator launched
        1995 : JavaScript created (10 days!)
             : PHP released
        1998 : Google founded
    section 2000s
        2004 : Gmail & Facebook launch
        2005 : Ajax popularised
             : YouTube founded
        2008 : Chrome browser released
             : V8 JS engine
        2009 : Node.js created
    section 2010s
        2013 : React released by Facebook
        2014 : Vue.js created
        2016 : Angular 2 released
    section 2020s
        2022 : Bun runtime released
        2023 : LLMs enter web tooling
        2025 : WebAssembly mainstream
\end{mermaid}

%% ============================================================
\section{Pie Chart}
%% ============================================================

\begin{mermaid}[type=pie, theme=sunset,
                caption={Traffic source breakdown},
                label={fig:pie}]
pie showData
    title Website Traffic Sources
    "Organic Search"  : 42.5
    "Direct"          : 21.0
    "Social Media"    : 17.3
    "Referral"        : 11.8
    "Email"           :  5.2
    "Paid Ads"        :  2.2
\end{mermaid}

%% ============================================================
\section{Mindmap}
%% ============================================================

\begin{mermaid}[type=mindmap, theme=forest,
                caption={Machine learning concept map},
                label={fig:mindmap}]
mindmap
  root((Machine Learning))
    Supervised
      Classification
        Decision Trees
        SVM
        Neural Nets
      Regression
        Linear
        Polynomial
        Ridge / Lasso
    Unsupervised
      Clustering
        K-Means
        DBSCAN
      Dimensionality Reduction
        PCA
        t-SNE
        UMAP
    Reinforcement
      Q-Learning
      Policy Gradient
      Actor-Critic
    Deep Learning
      CNNs
      RNNs / LSTMs
      Transformers
      Diffusion Models
\end{mermaid}

%% ============================================================
\section{Git Graph}
%% ============================================================

\begin{mermaid}[type=git, theme=arctic,
                caption={Feature branch git workflow},
                label={fig:git}]
gitGraph
    commit id: "Initial commit"
    commit id: "Add CI pipeline"

    branch develop
    checkout develop
    commit id: "Setup project"
    commit id: "Add auth module"

    branch feature/payments
    checkout feature/payments
    commit id: "Stripe integration"
    commit id: "Add webhooks"
    commit id: "Tests pass"

    checkout develop
    merge feature/payments id: "Merge payments"

    branch feature/dashboard
    checkout feature/dashboard
    commit id: "Dashboard UI"
    commit id: "Charts added"

    checkout develop
    merge feature/dashboard id: "Merge dashboard"
    commit id: "Integration tests"

    checkout main
    merge develop id: "Release v1.2.0" tag: "v1.2.0"
    commit id: "Hotfix: null check" type: HIGHLIGHT
\end{mermaid}

%% ============================================================
\section{Themes Showcase}
%% ============================================================

The five built-in themes can be selected per-diagram using
the \texttt{theme=\{name\}} option or globally with
\verb|\usepackage[theme=name]{mermaid-overleaf}|.

\subsection{All Five Themes — Same Diagram}

\begin{mermaid}[type=flowchart, theme=ocean, caption={Theme: ocean (default)}]
flowchart LR
    A[Input] --> B{Valid?} -->|Yes| C[Process] --> D[Output]
    B -->|No| E[Error]
\end{mermaid}

\begin{mermaid}[type=flowchart, theme=forest, caption={Theme: forest}]
flowchart LR
    A[Input] --> B{Valid?} -->|Yes| C[Process] --> D[Output]
    B -->|No| E[Error]
\end{mermaid}

\begin{mermaid}[type=flowchart, theme=sunset, caption={Theme: sunset}]
flowchart LR
    A[Input] --> B{Valid?} -->|Yes| C[Process] --> D[Output]
    B -->|No| E[Error]
\end{mermaid}

\begin{mermaid}[type=flowchart, theme=arctic, caption={Theme: arctic (light)}]
flowchart LR
    A[Input] --> B{Valid?} -->|Yes| C[Process] --> D[Output]
    B -->|No| E[Error]
\end{mermaid}

\begin{mermaid}[type=flowchart, theme=midnight, caption={Theme: midnight}]
flowchart LR
    A[Input] --> B{Valid?} -->|Yes| C[Process] --> D[Output]
    B -->|No| E[Error]
\end{mermaid}

%% ============================================================
\section{Using Pre-Rendered Images}
%% ============================================================

When you need publication-quality output, render your diagram
at \href{https://mermaid.live}{mermaid.live}, export as PDF or
SVG, upload to Overleaf, then use the \texttt{image} option:

\begin{verbatim}
\begin{mermaid}[
    image=figures/my-diagram.pdf,
    type=flowchart,
    theme=ocean,
    caption={My rendered diagram},
    label={fig:rendered},
    width=0.8\linewidth
]
flowchart TD
    A --> B --> C
\end{mermaid}
\end{verbatim}

The source listing is still shown below the image for
documentation purposes. Remove \BODY{} content to hide it.

\MermaidLiveNote

%% ============================================================
\section{Inline Code}
%% ============================================================

Use \verb|\mermaidcode{...}| to typeset Mermaid keywords inline:
in a flowchart, \mermaidcode{graph TD} sets top-down direction,
while \mermaidcode{graph LR} gives left-to-right layout.
Arrow types include \mermaidcode{-->} (solid),
\mermaidcode{-.->} (dotted), and \mermaidcode{==>} (thick).

%% ============================================================
\section{Quick Reference}
%% ============================================================

\begin{table}[H]
\centering
\caption{Environment options summary}
\begin{tabular}{lll}
\hline
\textbf{Option} & \textbf{Values} & \textbf{Description} \\
\hline
\texttt{type}    & \texttt{flowchart}, \texttt{sequence}, \texttt{gantt}, \ldots & Diagram type badge \\
\texttt{theme}   & \texttt{ocean}, \texttt{forest}, \texttt{sunset}, \texttt{arctic}, \texttt{midnight} & Colour theme \\
\texttt{image}   & file path & Pre-rendered image to include \\
\texttt{caption} & text & Figure caption \\
\texttt{label}   & text & \texttt{\textbackslash label} for \texttt{\textbackslash ref} \\
\texttt{width}   & dimension (default \texttt{0.85\textbackslash linewidth}) & Listing/image width \\
\texttt{pos}     & \texttt{H}, \texttt{h}, \texttt{t}, \texttt{b} & Float placement \\
\hline
\end{tabular}
\end{table}

\end{document}
